% Part: second-order-logic
% Chapter: sol-and-set-theory
% Section: cardinalities

\documentclass[../../../include/open-logic-section]{subfiles}

\begin{document}

\olfileid{sol}{set}{crd}

\olsection{توان‌های مجموعه‌ها}

\begin{explain}
همان‌گونه که می‌توانیم متناهی یا نامتناهی و !!{enumerable} یا
!!{nonenumerable}بودن دامنه را بیان کنیم، می‌توانیم ویژگی متناهی یا
نامتناهی و !!{enumerable} یا !!{nonenumerable}بودن زیرمجموعه‌ای از
$\Domain{M}$ را تعریف کنیم.
\end{explain}

\begin{prop}
فرمول~$\fn{Inf}(X) \ident$
\begin{multline*}
\lexists[u][(\lforall[x][\lforall[y][(\eq[u(x)][u(y)] \lif
      \eq[x][y])]] \land {} \\ \lforall[x][(X(x) \lif X(u(x)))]
      \land {} \\ \lexists[y][(X(y) \land \lforall[x][(X(x)
      \lif \eq/[y][u(x)])])])]
\end{multline*}
نسبت به تخصیص متغیر~$s$ ارضا می‌شود اگر و تنها اگر~$s(X)$ نامتناهی
باشد.
\end{prop}

\begin{prop}
فرمول~$\fn{Count}(X) \ident $
\begin{multline*}
(\lnot \lexists[x][X(x)] \lor \lexists[z][\lexists[u][(X(z) \land
    \lforall[x][(X(x) \lif X(u(x)))] \land {} \\ \lforall[Y][((Y(z) \land
      \lforall[x][(Y(x) \lif Y(u(x)))]) \lif X \subseteq Y)])]])
\end{multline*}
نسبت به تخصیص متغیر~$s$ ارضا می‌شود اگر و تنها اگر~$s(X)$
!!{enumerable} باشد.
\end{prop}

از قضیهٔ کانتور می‌دانیم مجموعه‌های !!{nonenumerable} وجود دارند و در
واقع بی‌نهایت سطح متفاوت از اندازه‌های نامتناهی در کار است. نظریهٔ
مجموعه‌ها حسابی کامل برای اندازهٔ مجموعه‌ها می‌پرورد و اعداد اصلی
نامتناهی را به مجموعه‌ها نسبت می‌دهد. اعداد طبیعی نقش اعداد اصلیِ
اندازه‌گیرِ مجموعه‌های متناهی را دارند. توان مجموعه‌های !!{denumerable}
نخستین توان نامتناهی است که~$\aleph_0$ («الفِ صفر») نامیده می‌شود. اندازهٔ
نامتناهی بعدی~$\aleph_1$ است. این کوچک‌ترین اندازه‌ای است که مجموعه‌ای
می‌تواند داشته باشد بی‌آنکه شمارا باشد (یعنی اندازهٔ~$\aleph_0$ داشته
باشد). می‌توانیم «$X$ اندازهٔ~$\aleph_0$ دارد» را با
$\fn{Aleph}_0(X) \liff \fn{Inf}(X) \land \fn{Count}(X)$ تعریف کنیم.
$X$ اندازهٔ~$\aleph_1$ دارد اگر و تنها اگر نامتناهی باشد، همهٔ
  زیرمجموعه‌هایش متناهی یا دارای اندازهٔ~$\aleph_0$ یا هم‌شمار با خود آن
  باشند، و خود اندازهٔ~$\aleph_0$ نداشته باشد. پس می‌توانیم این امر را
با فرمول~$\fn{Aleph_1}(X) \ident \fn{Inf}(X) \land
\lforall[Y][(Y \subseteq X \lif (\lnot\fn{Inf}(Y) \lor
\fn{Aleph}_0(Y) \lor \cardeq{X}{Y}))] \land \lnot
\fn{Aleph}_0(X)$ بیان کنیم. اندازهٔ~$\aleph_2$ نیز به‌شیوه‌ای مشابه
تعریف می‌شود، و الی آخر.

یک اندازه اهمیت ویژه دارد که به آن توان پیوستار می‌گویند. این اندازهٔ
$\Pow{\Nat}$ یا، به‌طور هم‌ارز، اندازهٔ~$\Real$ است. اندازهٔ
پیوستارداشتن یک مجموعه را نیز می‌توان در منطق مرتبهٔ دوم بیان کرد، اما
این کار اندکی بیشتر زحمت می‌طلبد.

\end{document}
